Fix an integer (n\ge 1).

A **feasible sequence** for (n) is a finite strictly increasing sequence of integers
[
1\le a_1<a_2<\cdots<a_k\le n
]
such that **all consecutive (interval) sums are distinct**, i.e. for every two *different* index pairs ((u,v)\neq(u',v')) with (1\le u\le v\le k) and (1\le u'\le v'\le k),
[
\sum_{i=u}^{v} a_i ;\neq; \sum_{i=u'}^{v'} a_i.
]

Define
[
f(n)=\max{k:\ \exists\ \text{a feasible sequence }(a_1,\dots,a_k)\text{ for }n}.
]

The Erdős–Harzheim problem (as stated on the site) asks for the asymptotic growth of (f(n)) as (n\to\infty), in particular:

* Does (f(n)=o(n)) (equivalently (f(n)/n\to 0))?

(There is also a variant (g(n)) where one drops the monotonicity assumption, but I treat it only as context.)

**Stress points / edge cases:**

* For (k=1) the condition is vacuous; so (f(n)\ge 1) always.
* For a given (k), there are (\binom{k+1}{2}=k(k+1)/2) consecutive sums to be all distinct.
* The condition is conveniently rephrased via prefix sums (see Lemma 1 below).
